\uuid{88Eh}
\titre{Constante de normalisation d'une densité gaussienne}
\auteur{Maxime Nguyen}
\organisation{AMSCC}
\datecreate{2026-11-27}
\niveau{L2}
\module{Probabilité}
\chapitre{Probabilité continue}
\sousChapitre{Loi normale}
\difficulte{2}
\theme{loi normale, densité, constante de normalisation}

\contenu{
\begin{question}
\qtype{onlyone}
On note $\mathcal{N}(\mu,\sigma)$ la loi normale d'espérance $\mu$ et d'écart-type $\sigma$ ; sa densité est $x \mapsto \frac{1}{\sigma\sqrt{2\pi}}\,e^{-\frac{(x-\mu)^2}{2\sigma^2}}$. Pour quelle valeur de la constante $c > 0$ la fonction $f(x) = c\,e^{-x^2}$, définie sur $\mathbb{R}$, est-elle une densité de probabilité ?
\begin{answers}
\bad{$c = \dfrac{1}{\sqrt{2\pi}}$}
\good{$c = \dfrac{1}{\sqrt{\pi}}$}
\bad{$c = \sqrt{\pi}$}
\bad{Aucune : $e^{-x^2}$ n'a pas la forme de la densité d'une loi normale.}
\end{answers}
\begin{explanations}
On identifie $e^{-x^2} = e^{-\frac{x^2}{2\sigma^2}}$
avec $2\sigma^2 = 1$, soit $\sigma = \frac{1}{\sqrt2}$ et $\mu = 0$. La
constante de $\mathcal{N}\left(0, \frac{1}{\sqrt2}\right)$ est
$\frac{1}{\sigma\sqrt{2\pi}} = \frac{\sqrt2}{\sqrt{2\pi}} = \frac{1}{\sqrt\pi}$.
On en déduit au passage $\int_{-\infty}^{+\infty} e^{-x^2}\,\mathrm{d}x = \sqrt\pi$.

A recopie la constante de $\mathcal{N}(0,1)$ sans tenir compte de $\sigma$.

C est la valeur de l'intégrale, et non son inverse.

D : toute fonction $x \mapsto c\,e^{-a(x-b)^2}$ avec $a > 0$ est, à la constante près, une densité normale.

Contre-sens visé : ne savoir reconnaître la loi normale que sous la forme
exacte du cours, sans identifier $\sigma$ dans l'exposant.
\end{explanations}
\end{question}
}
